\documentclass[../Main.tex]{subfiles}
\begin{document}

\section{Motivation}
\label{section:1.1}

In recent years, Vietnam has faced increasingly severe and unpredictable natural disasters, particularly in the central region. The year 2025 was recorded as one of the most extreme years for meteorological anomalies, breaking numerous historical records. According to the Ministry of Agriculture and Rural Development, natural disasters in 2025 caused estimated economic damages of up to 97,000 billion VND, resulting in 419 casualties and missing persons. Extreme weather events included 21 storms and tropical depressions in the East Sea, alongside unprecedented rainfall, such as the 1,739mm recorded within 24 hours at the Bach Ma station. These extreme floods cause widespread devastation, frequently destroying local telecommunications infrastructure and leading to prolonged power outages. Consequently, trapped victims are completely isolated in environments with intermittent or entirely disconnected Internet access. The inability to transmit distress signals or pinpoint exact locations severely delays rescue operations, directly threatening the survival of stranded individuals. Establishing a resilient communication method to track and locate victims in these harsh conditions is a critical necessity.

Alongside the immediate need for rescue operations, the post-disaster relief and charity mobilization process presents another significant societal challenge regarding financial transparency. Recently, the mobilization of relief funds by high-profile individuals and voluntary organizations has sparked widespread public skepticism. Instances where tens to hundreds of billions of VND were collected without independent auditing, detailed bank statements, or clear transaction records have led to profound crises of public trust. Furthermore, the lack of a standardized system allows for the unverified distribution of goods, while mandatory but untracked micro-donations in educational institutions highlight the systemic absence of an accountable framework. The current manual and fragmented approaches to managing charity campaigns fail to provide traceability. Therefore, developing a centralized platform to strictly manage the lifecycle of relief funds—from donation to final distribution—is essential to restore community trust and ensure resources reach the intended beneficiaries transparently and legally.

\section{Objectives and scope of the graduation thesis}
\label{section:1.2}

Currently, technological solutions supporting flood rescue and relief management in Vietnam exhibit critical technical and functional limitations. Instant messaging applications like Zalo require active user intervention to share locations within a chat interface, which becomes impossible during network outages and lacks direct integration with official rescue teams. Dedicated emergency platforms, such as web-based SOS systems, demand a stable Internet connection to trigger alerts and immediately fail when the network is disrupted. They also lack continuous background tracking and family-sharing functionalities. Social location-tracking applications, such as "Whoo", offer features to save the last known location and share it among friends; however, they suffer from system instability and do not provide an interface for professional rescue coordination. Crucially, none of these existing platforms incorporate features to manage and monitor the subsequent charity and relief distribution processes, leaving a significant gap in the disaster response lifecycle.

To overcome these limitations, this thesis aims to design and develop FloodHelper, a comprehensive ecosystem combining offline-resilient rescue coordination and transparent charity management. The system's scope encompasses three main interfaces tailored to specific user roles. First, a mobile application is developed for Victims, Rescuers, and Benefactors, facilitating location broadcasting and campaign interactions. Second, a web-based portal is designed specifically for Authorities to verify user roles, approve rescue requests, and monitor operations. The administrative dashboard for system-wide user management is considered as future work.

In terms of business logic, the charity management subsystem strictly enforces a transparent workflow. A Benefactor creates a campaign, which remains inactive until officially approved by an Authority. Once the donation phase begins, all financial contributions are publicly recorded. Subsequently, the allocation phase requires the Benefactor to publicly declare the distribution of funds and specific relief items. The campaign concludes only when legitimate bank statements are uploaded for final verification. To ensure absolute integrity, Authorities are granted the power to suspend any campaign exhibiting suspicious activities during any phase.

\section{Tentative solution}
\label{section:1.3}

To address the challenges of intermittent connectivity during floods, the proposed solution utilizes the Message Queuing Telemetry Transport (MQTT) protocol as the core communication foundation. MQTT is selected due to its lightweight payload, low bandwidth requirements, and robust message retention capabilities, making it exceptionally reliable in weak network conditions. On the client side, the mobile application leverages background isolates to maintain persistent background processes, periodically capturing and pushing the user's GPS coordinates to the MQTT broker without requiring active screen time. All critical actions, including broadcasting SOS signals, handling rescue assignments, and terminating alerts, are executed via MQTT topics, which subsequently synchronize with the backend server whenever the network connection is briefly restored.

For the charity management module, the backend architecture implements automated scheduling mechanisms, utilizing cron jobs to securely manage campaign state transitions (e.g., from "Approved" to "Donating" and "Distributing") based on predefined temporal constraints and verification rules, minimizing manual interference.

The primary contributions of this graduation thesis are threefold: (i) designing a resilient mobile architecture utilizing MQTT and background isolates to ensure continuous location tracking and SOS broadcasting in unstable network environments; (ii) digitizing the traditional charity process into a strictly state-driven, publicly transparent workflow that enforces accountability from the initial donation to the final bank statement validation; and (iii) successfully integrating these two critical components into a unified, functional platform (FloodHelper) to comprehensively support both immediate disaster rescue and post-disaster relief efforts.

\section{Thesis organization}
\label{section:1.4}

The remainder of this graduation thesis report is organized as follows.

Chapter 2 presents the requirement survey and analysis phase of the FloodHelper system. This chapter focuses on examining the practical needs of victims in flooded areas and evaluating the limitations of current rescue applications, thereby establishing detailed functional and non-functional requirements for both the rescue and charity management subsystems.

In Chapter 3, I introduce the methodology applied to resolve the identified challenges. The chapter details the selection of the MQTT protocol for communication in weak network environments and outlines the technological approaches utilized to build a robust architecture for data synchronization and transparent charity tracking.

Chapter 4 details the system design, implementation process, and practical evaluation. This chapter encompasses the overall architectural diagrams, database modeling, and application development, followed by testing results that assess the system's communication performance under simulated intermittent network conditions.

Chapter 5 summarizes the finalized solutions and the primary contributions of the thesis. The content highlights the successful integration of the resilient location-sharing mechanism and the comprehensive charity management dashboard, demonstrating how these components directly address the initial research problems.

Finally, Chapter 6 provides the conclusion of the study and outlines future work. This chapter reflects on the achieved results, acknowledges the current limitations of the system, and proposes potential research directions, including the integration of Delay-Tolerant Networking (DTN) to support data sharing in completely disconnected scenarios.

\end{document}